\documentclass[../../main.tex]{subfiles}
\begin{document}

\begin{center}
  {\fontsize{48pt}{58pt}\selectfont \textbf{2010}\normalsize\textbf{年度}}
\end{center}

\begin{table}[htb]
  \centering
  \caption{出題テーマと難易度}
  \begin{tabular}{|c|p{9cm}|c|}
    \hline
    \textbf{問題番号} & \textbf{テーマ} & \textbf{難易度} \\
    \hline
    第1問 & 【甲・乙共通】四面体と直交平面／【甲】数の並べ方 & \\
    \hline
    第2問 & 三角形の角の最大値 & \\
    \hline
    第3問 & 正弦曲線と余弦曲線で囲む面積 & \\
    \hline
    第4問 & 【乙】鋭角三角形の辺長／【甲】数列の不等式 & \\
    \hline
    第5問 & 三の冪の二進付値 & \\
    \hline
    第6問 & 【乙】箱への球の配置確率／【甲】立方体の回転体 & \\
    \hline
  \end{tabular}
\end{table}

\begin{table}[htb]
  \centering
  \caption{解答}
  \renewcommand{\arraystretch}{1.4}
  \begin{tabular}{|c|c|p{9cm}|}
    \hline
    \textbf{問題番号} & \textbf{小問} & \textbf{答え} \\
    \hline
    第1問 & 【甲・乙共通】 & 平面$ABM$は辺$CD$に垂直（証明問題） \\
    \cline{2-3}
    & 【甲】 &  \\
    \hline
    第2問 & -- & $\max\angle APB=\dfrac{\pi}{4}$ \\
    \hline
    第3問 & -- & $a=\dfrac{4}{3}$ \\
    \hline
    第4問 & 【乙】 & $b=\dfrac{a+\sqrt{3(4-a^2)}}{2}$ \\
    \cline{2-3}
    & 【甲】 &  \\
    \hline
    第5問 & 【乙】(1) & $v_2(3^{2^n}-1)=n+2$（証明問題） \\
    \cline{2-3}
    & 【乙】(2) & $m=2,4$ \\
    \hline
    第6問 & 【乙】 & $\lim_{n\to\infty}\dfrac{\log p_n}{n}=\log2-1$ \\
    \cline{2-3}
    & 【甲】 &  \\
    \hline
  \end{tabular}
\end{table}

\paragraph{試験形態}
6問．試験時間は要確認．

\end{document}
